\documentclass[11pt]{article}
\usepackage[margin=1in]{geometry}
\usepackage{amsmath,amssymb,amsthm}
\usepackage[hidelinks]{hyperref}
\newtheorem{theorem}{Theorem}
\newtheorem{lemma}{Lemma}
\newcommand{\J}{\mathcal J}
\title{Counterexamples to the proposed Redheffer lower bound\\
for a ratio of normalized Bessel functions}
\author{Principia Math In Progress}
\date{18 September 2026}
\begin{document}
\maketitle
\begin{abstract}
The lower bound in MathDB 335189 fails at the explicit pair
\((\nu,x)=(30,1)\). We certify this with rational arithmetic and
a variational bound for the first Bessel zero. More generally, for
every real \(\nu\ge28\) it fails at all sufficiently small nonzero
real \(x\). The upper bound was already known in the literature.
The mathematical counterexample is complete, but campaign eligibility
remains pending inspection of the complete 2018 journal version of
the source. No new research credit or priority is claimed here.
\end{abstract}

\section*{Task 1: the conjecture and a complete counterexample}
For \(\nu>-1\), define the even entire normalized Bessel function by
\[
 \J_\nu(x)=\sum_{k=0}^{\infty}
 \frac{(-1)^k x^{2k}}{4^k k!(\nu+1)_k},
 \qquad
 \J_\nu(x)=2^\nu\Gamma(\nu+1)x^{-\nu}J_\nu(x)\quad(x>0).
\]
Here \((a)_k=a(a+1)\cdots(a+k-1)\), with \((a)_0=1\).
Write \(j=j_{\nu,1}\) for the first positive zero of \(J_\nu\),
\(q=j^2\), and \(R_\nu(x)=\J_{\nu+1}(x)/\J_\nu(x)\).
Conjecture 1, equation (3.2), on page 8 of the current preprint
\cite{BM} asserts, for \(|x|<j\),
\begin{equation}\label{conj}
 \left(\frac q{q-x^2}\right)^{q/[4(\nu+1)(\nu+2)]}
 \le R_\nu(x)\le\frac q{q-x^2}.
\end{equation}
The source explicitly conjectures the left inequality for every
\(\nu>-1\); our examples satisfy that intended range.

\begin{theorem}\label{main}
For every real \(\nu\ge28\), there is \(\varepsilon_\nu>0\)
such that the left inequality in \eqref{conj} is strictly reversed
whenever \(0<|x|<\varepsilon_\nu\).
It is also strictly reversed at \(\nu=30,\ x=1\).
\end{theorem}

\subsection*{A bound for the first zero}
\begin{lemma}\label{zero}
For every real \(\nu\ge28\),
\[
 j_{\nu,1}^2\le U_\nu:=\frac{3(2\nu+7)(2\nu+8)}{10}
 <T_\nu:=\frac{4(\nu+1)(\nu+2)(\nu+3)}{3\nu+5}.
\]
\end{lemma}
\begin{proof}
Set \(\phi(r)=J_\nu(jr)\), which is positive on \(0<r<1\).
The Bessel equation gives
\[
 -(r\phi')'+\frac{\nu^2}{r}\phi=j^2r\phi .
\]
For \(u(r)=r^\nu(1-r)^3\), integration by parts gives the
ground-state identity
\begin{equation}\label{ground}
 \int_0^1\left(r(u')^2+\frac{\nu^2}{r}u^2-j^2ru^2\right)\,dr
 =\int_0^1r\phi^2\left((u/\phi)'\right)^2\,dr\ge0.
\end{equation}
For completeness, substitute \(u=\phi v\) in the left side:
its integrand is \(r\phi^2(v')^2+(r\phi\phi'v^2)'\).
The boundary term \(r\phi'u^2/\phi\) tends to zero at both ends.
At zero, \(\phi(r)=c r^\nu(1+O(r^2))\); at one,
\(\phi(r)\) has a simple zero and \(u(r)=O((1-r)^3)\).
All the displayed integrals converge. Simplicity of the zero follows
from uniqueness for the Bessel differential equation at \(r=1\).

Writing \(u=r^\nu w\), integration of the cross term yields
\[
 \int_0^1\left(r(u')^2+\frac{\nu^2}{r}u^2\right)\,dr
 =\int_0^1r^{2\nu+1}(w')^2\,dr,
 \qquad w=(1-r)^3.
\]
Indeed, the removed boundary term is \(\nu r^{2\nu}w^2\);
it vanishes at both endpoints. Consequently \eqref{ground} implies
\[
 j^2\le
 \frac{9\int_0^1r^{2\nu+1}(1-r)^4\,dr}
      {\int_0^1r^{2\nu+1}(1-r)^6\,dr}
 =9\frac{B(2\nu+2,5)}{B(2\nu+2,7)}
 =U_\nu.
\]
The beta-integral evaluation follows by repeated integration by parts.
Finally, if \(s=\nu-28\ge0\), direct algebra gives
\[
 10(3\nu+5)(T_\nu-U_\nu)
 =4s^3+246s^2+3854s+2256>0.
\]
\end{proof}

\subsection*{Failure near zero for every order at least 28}
Both normalized functions are positive near zero, so their real
logarithms have convergent Taylor series there. The defining series gives
\[
 \log R_\nu(x)=
 \frac{x^2}{4(\nu+1)(\nu+2)}
 +\frac{(3\nu+5)x^4}
 {32(\nu+1)^2(\nu+2)^2(\nu+3)}+O(x^6).
\]
The logarithm of the proposed lower bound is
\[
 \frac{q}{4(\nu+1)(\nu+2)}
       \log\frac{q}{q-x^2}
 =\frac{x^2}{4(\nu+1)(\nu+2)}
  +\frac{x^4}{8(\nu+1)(\nu+2)q}+O(x^6).
\]
Thus their difference has coefficient
\[
 C_\nu=
 \frac{(3\nu+5)q-4(\nu+1)(\nu+2)(\nu+3)}
 {32(\nu+1)^2(\nu+2)^2(\nu+3)q}
\]
at \(x^4\). Lemma \ref{zero} proves \(C_\nu<0\).
For each fixed \(\nu\ge28\), analyticity therefore makes the difference
strictly negative at every sufficiently small nonzero real \(x\).
We can choose \(\varepsilon_\nu<j_{\nu,1}\), proving the first assertion.
This is an analytic sign argument, not an inference from a finite
set of Taylor coefficients or numerical samples.

\subsection*{An exact counterexample at \((30,1)\)}
Set \(q=j_{30,1}^2\), \(U=6834/5\), and \(A=3968=4\cdot31\cdot32\).
Lemma \ref{zero} gives \(q\le U\). Also \(j_{30,1}>1\):
for \(0\le x\le1\), the alternating defining series has decreasing
term magnitudes and
\(\J_{30}(x)\ge1-x^2/124>0\).
The claimed inequality is therefore defined at \(x=1\).

Let
\[
 S_{\nu,m}=\sum_{k=0}^m\frac{(-1)^k}{4^k k!(\nu+1)_k}.
\]
The alternating-series remainder bound implies
\[
 0<S_{30,5}\le\J_{30}(1),\qquad
 \J_{31}(1)\le S_{31,4}.
\]
The exact partial sums are
\[
 S_{30,5}=\frac{4748440055099}{4786893619200},
 \qquad S_{31,4}=\frac{7660705361}{7720796160}.
\]
Consequently, with \(d=252139378/10^{12}>0\),
\[
 R_{30}(1)\le
 \frac{4749637323820}{4748440055099}<1+d.
\]
One exact check of the strict inequality is
\[
 d-\left(\frac{4749637323820}{4748440055099}-1\right)
 =\frac{481473794211}{2374220027549500000000000}>0.
\]
For \(d\ge0\), integration of
\((1+t)^{-1}\le1-t+t^2\) gives
\(\log(1+d)\le d-d^2/2+d^3/3\). On the other hand, since
\(1<q\le U\),
\[
 \frac qA\log\frac q{q-1}
 =\frac1A\sum_{m=1}^\infty\frac{q^{1-m}}m
 \ge\frac1A+\frac1{2AU}.
\]
The remaining rational comparison is
\[
 \frac1A+\frac1{2AU}
       -\left(d-\frac{d^2}{2}+\frac{d^3}{3}\right)
 =\frac{9597654965608638275835585186379}
 {13240875000000000000000000000000000000000}
 >\frac7{10^{10}}>0.
\]
It follows that
\[
 \log R_{30}(1)<\frac qA\log\frac q{q-1},
\]
which proves the explicit counterexample and completes the proof
of Theorem \ref{main}. The accompanying verifier reproduces every
rational comparison and the symbolic identities.

\subsection*{Scope of the conclusion and the known upper bound}
A single admissible counterexample disproves \eqref{conj} as a universal
statement; the interval family supplies a second analytic proof.
We do not assert failure at every order below 28.
For context, the upper bound in \eqref{conj} is already credited
to earlier work by \cite{BM}, page 8.
Its short proof uses the classical Mittag--Leffler expansion
and first Rayleigh sum recorded in that source:
\[
 R_\nu(x)=\sum_{k\ge1}
 \frac{w_k}{1-x^2/j_{\nu,k}^2},\qquad
 w_k=\frac{4(\nu+1)}{j_{\nu,k}^2}>0,\qquad \sum_{k\ge1}w_k=1.
\]
Every denominator is at least \(1-x^2/j_{\nu,1}^2\) for real
\(|x|<j_{\nu,1}\), proving the upper inequality. Its exponent 1
is necessary because the first summand has a simple pole at
\(x=j_{\nu,1}\). This known upper bound earns no separate credit.

\section*{Task 2: determine the exact range of valid orders}
Find the set of \(\nu>-1\) for which the proposed lower bound
holds throughout \((-j_{\nu,1},j_{\nu,1})\).
The inequality \(j_{\nu,1}^2\ge T_\nu\) is a necessary local condition.
Determine whether it is sufficient and, if not, identify the
additional obstruction. The proved \(\nu=-1/2\) case in \cite{BM}
provides a starting point.

\section*{Task 3: determine the true optimal lower exponent}
Study
\[
 a_\nu=\inf_{0<x<j_{\nu,1}}
 \frac{\log R_\nu(x)}
      {\log(j_{\nu,1}^2/(j_{\nu,1}^2-x^2))}.
\]
Determine the minimizer and a useful formula or certified bounds.
For \(\nu\ge28\), the proof shows that \(a_\nu\) is strictly less
than the proposed value \(j_{\nu,1}^2/[4(\nu+1)(\nu+2)]\).
Resolve its behavior as \(\nu\) tends to infinity.

\section*{Task 4: classify the Rayleigh-coefficient monotonicity}
The source proposes monotonicity in \(m\) of
\[
 \omega_{m,\nu}=j_{\nu,1}^{2m}
   \left(\sum_{k\ge1}j_{\nu,k}^{-2m}
               -\sum_{k\ge1}j_{\nu+1,k}^{-2m}\right).
\]
The inequality \(\omega_{2,\nu}\ge\omega_{1,\nu}\) is exactly the
necessary condition \(j_{\nu,1}^2\ge T_\nu\).
Classify the orders and indices where monotonicity changes, and
determine whether the sequence has at most one change of direction.

\section*{Task 5: optimize elementary variational certificates}
Replace \(r^\nu(1-r)^3\) by \(r^\nu(1-r)^p\) with \(p>1/2\).
The same computation gives
\[
 j_{\nu,1}^2\le
 \frac{p(2\nu+2p+1)(2\nu+2p+2)}{2(2p-1)}.
\]
Optimize this bound, or extend the trial space by a polynomial
factor, to obtain the smallest provable threshold for counterexamples
using elementary exact inequalities.

\begin{thebibliography}{9}
\bibitem{BM}
\'A.~Baricz and K.~Mehrez,
\emph{Redheffer type bounds for Bessel and modified Bessel functions
of the first kind}, arXiv:1701.08446v1 (2017), 9 pp.,
\url{https://arxiv.org/abs/1701.08446}.
Complete preprint inspected; Conjecture 1 is on p.~8.
Published in \emph{Aequationes Mathematicae} 92 (2018), 425--439,
\url{https://doi.org/10.1007/s00010-018-0545-8}.
The complete journal text has not been inspected.
\end{thebibliography}
\end{document}
